\chapter{Mitschriften von Mitarbeiteraussagen während der Befragung und des Usability-Tests}
\label{Add:Mitschriften_Mitarbeiteraussagen}

% https://tex.stackexchange.com/a/39645
\lstset{literate=%
    {Ö}{{\"O}}1
    {Ä}{{\"A}}1
    {Ü}{{\"U}}1
    {ß}{{\ss}}1
    {ü}{{\"u}}1
    {ä}{{\"a}}1
    {ö}{{\"o}}1
}

\begin{lstlisting}[language={[LaTeX]TeX},inputencoding={utf8}]
Zusammenfassung: Feedback / Wünsche für das Dashboard:
- Security Error beim öffnen von Grafana
- die letzte Spalte [test_cat_5] kann statt des Kürzels Namen [...] bekommen - ID 4
- 2. Dashboard wurde oft nicht wahrgenommen / vergessen / nicht mehr gefunden später
- die ignorierten Werte können aus der Ansicht rausgenommen werden
- Layout verrutscht bei größeren Werten (Bsp.: "529", "crashed")
- Tree-verwöhnt - "warum werden die  anderen Dashboards nicht als tree angezeigt?"
  --> Breadcrumbs/Pfadangabe reichen nicht aus
- DB-Benennung (dev) aus dem alten System übernommen -> aktuelle Version!
- Darkmode - gibts Grafana im Lightmode?
- hellgrauen Kasten um die zusammengehörigen Zahlen einer Kategorie machen
- Datumsauswahl gibt es schon
  --> ABER: wie sich das Ergebnis von gestern auf heute geändert hat
      (nicht aber von gestern auf vorgestern)
  --> allgemein Verlaufsgrafiken der letzten 14 Tage
- Grafiken einbauen, statt nur Zahlen
  --> erst ab Verlaufsdaten möglich

=======

Zusammenfassung: Green Day (Aufgabe 1) Begründungen:
- nein
  --> es sind nicht alle Tests grün / 0 / fehlerfrei
- ja
  --> bekannte Fehler - Bsp. Timingproblem in der Testumgebung
  --> TUI Umgebung wird nicht beachtet (sehr instabile Tests)
  ==> Wissen benötigt, was ignoriert werden darf - es sollte langfristig lieber die Testumgebung aufgeräumt werden

=======

Testperson 1:
-------------
altes System:
- nutzt das seit 30+ Jahren
- "ich mache mir jeden Morgen 20 [Oberflächen / Tabs] auf" - "komme in 3 Minuten durch"
  --> muss täglich viele der Ergebnisse durchschauen
  --> ist mit für die Wartung zuständig
- seit 2023 haben die Kollegen was noch nicht gefixt - Verschiebung
- "ich habe eine Vorgehensweise wie ich das alles lese"
  --> "fotografischer Blick"
Dashboard:
- "ich muss wirklich lesen, ich habe den Eindruck, dass ich sonst nicht mehr lese, wenn ich die [Oberflächen des alten System] durchgucke"
- "das ist jetzt in Spalten angeordnet, was vorher in Zeilen war"

Testperson 2:
-------------
- macht das noch länger als [Testperson 1]
altes System:
- Green Day: "nach den Regeln der Kunst nein, aber mit beiden Augen zudrücken ja"
- Anfang war nur die BUILD-Kategorie
  --> ist immer weiter gewachsen, man übersieht auch gerne mal eine Zahl
  --> für den groben Überblick gut geeignet
Dashboard:
- gut, irritiert hat: das Zeilen Spalten vertauscht sind
- deutlich einfacher, als das andere, ob noch Tests fehlen
- durch das simplifizieren der Ausgabe

Testperson 3:
-------------
altes System:
- ist ein Green Day mit Spezialwissen
- 30 Jahre mit [altes System] und Jenkins
- standardmäßig 5 Tabs (nächtlich laufend) - auf den ersten Blick bearbeitbar
- "viele alte Information und vieles führt ins Leere"
Dashboard:
- neues Dashboard natürlich erstmal gewöhnungsbedürftig
- Dashboard hat verwirrt, ist einiges an Information drauf
  --> viele Information die man vorher (im [altes System]) erstmal nicht beachtet/ausblendet
  --> man muss sich neu daran gewöhnen
- Zahlen sind sehr groß (man sieht zu wenig auf einmal)

Testperson 4:
-------------
altes System:
- Green Day - kann er nicht eindeutig erkennen
  --> nur [Testperson 3] weiß ob angezeigte Fehler in [test_cat_3] eine Relevanz haben
  --> genauso bei einem anderen System [kein Part des untersuchten Use-Cases]
  --> ohne das Wissen ist kein Green Day
- Maskenwechsel: "ich mache das mit den Tabs da geht das relativ schnell"
  --> gemeint: Tableiste ist links im Browser - Tabs untereinander aufgelistet
  --> Anmerkung Autorin: "ja aber sind viele Tabs und alle kryptisch benannt"
- hat für jede Version den Testpool abgespeichert
  --> auf VM, läuft 24/7
  --> ist alles gruppiert in Lesezeichen
Dashboard:
- "ist man gar nicht gewohnt so ein Comfort"
- wenn man in einem Dashboard ist sieht man die anderen nicht
  --> es gibt keine Tree-Funktonion, aber Breadcrumbs sind da
- schnell erlernbar, nicht vergleichbar
- wenn man sich 1-2 Tage einfindet, ist es super
- Wunsch: Overview über die einzelnen Versionen und deren aktuellen Zustand
  --> in der DB Liste: Overview OK oder nicht [wurde noch nicht geprüft ob technisch umsetzbar]

Testperson 5:
-------------
altes System:
- Green Day: sieht ok aus, nur die [test_cat_3] Tests - ja
Dashboard:
- findet gut dass man sich die Rechner nicht merken muss
- das Dashboard zu haben findet er gut
  --> Farbcodierung
  --> einheitliche UI/Navigieren
  --> filtert die unwichtigen Details aus
  --> anders sortiert, leichter zu verstehen
- "Darkmode mögen jüngere lieber"
- Grafiken einbauen, statt nur Zahlen
  --> erst ab Verlaufsdaten sinnvoll möglich - noch nicht implementiert

Testperson 6:
-------------
altes System:
- wenn rot ist das kein GreenDay
  --> er arbeitet aber auch nicht damit [mit explizit diesem System, aber dem Testpool], kann das nicht einschätzen
- eigentlich müsste man die Sachen kontinuierlich verbessern, als das es irgendwann heißt jetzt ein halbes Jahr um alles Umzukrempeln
  --> er hätte statt einem Executable einen REST mit React irgendwie geschrieben, was das von sich aus schon anders aufbaut
Dashboard:
- Beschreibung der Versionen ist besser
  --> statt Keys, oder Maschinen durchklickt
- Wunsch: soll man direkt drauf kommen von der Startseite, nicht erst über den Reiter
  --> Anmerkung Autorin: wird Grafana aufgerufen, muss sich (nach aktuellem Stand) zuerst angemeldet werden, dann den Reiter "Dashboards" auswählen und schließlich das richtige Dashboard aus der Liste (bestehend aus aktuell 2 DBs) anklicken
- man erkennt dass es ein Framework ist, mit dem man einiges mehr machen kann
  --> die Werbung ist unschön

Testperson 7:
-------------
altes System:
- Green Day: ja, aber ist eine Definitionsfrage
-> kommt auf die erfahrung an -> build, Batch und wichtige Tests waren ok
-> nur unwichtige und bekannte tests waren rot
neues System:
- findet gut:
  --> Farbgebung gut - Errors rot und pending so gelb-artig (Modi gut unterscheidbar)
  --> [altes System] muss man 2x klicken, deswegen DB besser
  --> aufklappen / reihen, findet er gut, geht im scresult nicht (war damals vor 20 Jahren Top-aktuell)
- Zahlen kleiner machen, dass man die 2 Modi nebeneinander / auf einen Blick sieht
  --> Anmerkung Autorin: ist via zoom machbar, dass 2 Reihen auf einmal sichtbar sind

Testperson 8:
-------------
- ist eher für [anderes System; kein Part des Use-Case] verantwortlich, kennt sich nur grob aus
altes System:
- Green Day: nein "Tests haben err"
  --> weil es rot leuchtet
- scresult Codes sind random, weiß man nicht
  --> hat die Links als Shortcuts und richtig benannt, macht das Dashboard besser
  --> Anmerkung Autorin: "Codes" = [altes System: Version, Maschine_1]
Dashboard:
- intuitivere Titel für Reihen, statt Codes (s. altes System)
- untersch. DBs waren nicht direkt ersichtlich, aber mit mehr Praxis, gewöhnt man sich dran

Testperson 9:
-------------
altes System:
- Green Day: viele Sachen laufen immer rot
  --> bedeutet nein, weil 4 Fehler / rote Anzeigen
    ----> GD = wenn die Seite green ist oder muss anders definiert sein
    ----> wenn was nicht wichtig wäre müsste es eigentlich rausgenommen werden
  --> "ich als Laie, würde das sagen"
- "es tendieren manche dazu, wenn ein Test länger rot ist, ihn zu ignorieren"
- esotherische Tests, die noch nie eine Rolle gespielt haben
  --> beim Kunden läuft es seit Jahrzehnten stabil, muss man pragmatisch sehen
  --> unterscheiden zwischen:
    ----> fehlerhaften Tests, die nicht einfach zu fixen sind
    ----> und den Altlasten, die seit Jahren nicht angefasst wurden und "ignoriert" werden können
- Umfrage: "das ist fies" [die bekannten / ignorierten Probleme so offensichtlich abzufragen]
  --> bekannte / offensichtliche / ignorierte Probleme
  --> neues Tool ist herausgehoben
Dashboard:
- DB sehr gut / übersichtlich, aber man muss sich einarbeiten
- Abkürzungen sind gewöhnungsbedürftig, wenn man nicht regelmäßig damit arbeitet
  --> [soll teilweise noch durch bessere Begriffe ausgetauscht werden]
- "ignored" evtl. auf neuen Tab, nimmt zu viel Platz weg für die Übersicht

Testperson 10:
--------------
altes System:
- kein Green Day 
  --> das ist rot, also muss Fehler

Testperson 11:
--------------
- es gibt 2 Maschinen
  --> nur das was auf der [Maschine 1] gebaut wird, wird an den Kunden ausgeliefert
  --> [Maschine 2] ist nur für die Tests
  --> Anmerkung Autorin: diese wurden hier als Modi bezeichnet, da zwei verschiedene Modi auf jeweils einer der Maschinen getestet wird
altes System:
- Green Day: unwahrscheinlich -> Fehler + CORE
  --> "ich müsste es als GreenDay bewerten" - auch wenn es das nicht ist
- Aufgabe 3: Kategorie gemeint - einer der Tests wurde angesehen
  --> wenn [Test] faul ist, braucht man die anderen nicht mehr angucken
  --> bauen auf dem ersten auf
- macht jeden Tag 14 Tabs auf
  --> befasst sich mehr mit den Windows-Tests, nicht so sehr mit dem Use-Case
  --> "das gefällt mir mit den Tabs gut, dass mache ich immer so"
- [Fehler]: ist seit 20 Jahren, kommt immer mal an einer anderen Stelle auf
  --> kommt fast jeden Tag
- jeder [Test-Prozess], muss seine Ergebnisse selber auf die Seite in Text schreiben
  --> braucht etwas zu lange, können sie sich die Ergebnisse gegenseitig wegschreiben
- altes System = results von source-control
  --> das ist ein ganz dummes Tool, das macht nicht viel
Dashboard:
- Übergang zwischen Dashboards: "finde ich schwierig: ein Dashboard ist offen, [das andere] gar nicht mehr im Blick gehabt"
- Namen brauchen viel Platz
- [Modus-Namen] nicht so geläufig - eher verwirrend gewesen
- kann die Hilfe nicht geben, die es bräuchte, um die Tests auszuwerten
  --> Unterschied zum Vortag
- 2 Fehler hast du immer
  --> nicht direkt erkennbar, ob relevant oder nicht [Anmerkung Autorin: dann sollte die Umgebung angepasst werden]
  --> Liste der deaktivierten Tests (Konzept funktioniert nicht, keiner geht die Liste durch um zu sehen, welche evtl. wieder funktionieren) - [wurden in "ignored" aufgelistet]
- Wunsch: 
  --> Datumsauswahl gibt es im alten System schon [im neuen noch nicht]
  --> ABER: wie sich das Ergebnis von gestern auf heute geändert habe (nicht aber von gestern auf vorgestern, das wäre zu dem Zeitpunk bereits bekannt)

Testperson 12:
--------------
- ist für einen ganz anderen Kontext zuständig (UI, Weboberfläche)
  --> Daten wurden aus allen Auswertungen rausgenommen
altes System:
- Green Day: nein, weil rot in der Tabelle sichtbar
- Aufgabe 3: "man kann nur raten - es ist nicht rot markiert"
--> könnte alles sein, müssen keine Fehler sein, die angezeigt werden
- es wäre was gutes, wenn man nicht voraussetzt, dass ein Fachfremder das nicht versteht
Dashboard:
- Dashboard ist falsch benannt
- Aufbereitung ist super, intuitiv, ist genau wie man sich das vorstellt

\end{lstlisting}
